\documentclass[../../../include/open-logic-section]{subfiles}

\begin{document}

\olfileid[id]{sth}{z}{power}
\olsection{Himpunan Kuasa}

Kita lanjutkan dengan aksioma lain:

\begin{axiom}[Himpunan Kuasa]
Untuk sebarang himpunan $A$, himpunan
$\Pow{A}=\Setabs{x}{x\subseteq A}$ ada.
\[
	\forall A \exists P \forall x(x \in P \liff (\forall z \in x)z \in A)
\]
\end{axiom}

Pembenaran bagi aksioma ini cukup langsung. Andaikan $A$ dibentuk pada
tahap~$S$. Maka semua anggota~$A$ telah tersedia sebelum~$S$ (berdasarkan
\stagesacc). Jadi, dengan penalaran seperti dalam pembenaran kita bagi
Pemisahan, setiap himpunan bagian dari~$A$ telah terbentuk paling lambat pada
tahap~$S$. Karena itu, semuanya tersedia untuk dibentuk menjadi satu himpunan
pada sebarang tahap sesudah~$S$. Kita tahu bahwa suatu tahap demikian ada,
sebab $S$ bukan tahap terakhir (berdasarkan \stagessucc). Jadi, $\Pow{A}$
ada.

Berikut suatu konsekuensi yang berguna dari Himpunan Kuasa:

\begin{prop}\label{thm:Products}
Dengan diberikan sebarang himpunan $A,B$, hasil kali Kartesius
$A\times B$ ada.
\end{prop}

\begin{proof}
Himpunan $\Pow{\Pow{A\cup B}}$ ada berdasarkan Himpunan Kuasa dan
\olref[sth][z][pairs]{prop:pairsconsequences}. Jadi, berdasarkan Pemisahan,
himpunan berikut ada:
\[
	C = \Setabs{z \in \Pow{\Pow{A \cup B}}}{(\exists x \in A)(\exists y \in B) z = \tuple{x, y}}.
\]
Untuk sebarang $x\in A$ dan $y\in B$, himpunan $\tuple{x,y}$ ada berdasarkan
\olref[sth][z][pairs]{prop:pairsconsequences}. Selain itu, karena
$x,y\in A\cup B$, kita mempunyai $\{x\},\{x,y\}\in\Pow{A\cup B}$, dan
$\tuple{x,y}\in\Pow{\Pow{A\cup B}}$. Jadi, $A\times B=C$.
\end{proof}

Dalam bukti ini, Himpunan Kuasa berinteraksi dengan Pemisahan. Hal itu tidak
mengherankan. Tanpa Pemisahan, Himpunan Kuasa bukanlah prinsip yang sangat
\emph{kuat}. Bagaimanapun, Pemisahan memberi tahu kita himpunan bagian mana
dari suatu himpunan yang ada, dan dengan demikian menentukan seberapa
``gemuk'' setiap himpunan kuasa.

\begin{prob}
Tunjukkan bahwa, untuk sebarang himpunan $A,B$: (i) himpunan semua relasi
dengan domain~$A$ dan daerah hasil~$B$ ada; dan (ii) himpunan semua fungsi
dari $A$ ke~$B$ ada.
\end{prob}

\begin{prob}
Misalkan $A$ suatu himpunan dan $\sim$ suatu relasi ekuivalensi pada~$A$.
Buktikan bahwa himpunan kelas-kelas ekuivalensi di bawah $\sim$ pada~$A$,
yakni $\equivclass{A}{\sim}$, ada.
\end{prob}

\end{document}
